\section{Counterfactual lesion audit}
\label{sec:audit}
% PLAN (about 1.5 pages). The definition of A is already in methods.tex (Lesion-absorption audit); this section states the protocol and its limits.
%   5.1 Definition and reading of A: slope of the change in the bone estimate on the inserted signal; A=0 keeps the lesion, A=1 removes it; A=0 is
%       reached by doing nothing, so A is always reported next to a suppression gain; A can leave [0,1].
%   5.2 Sites and lesions: lesions from held-out CTs (library built from the TEST split, s4/lesion_lib.py), six sites per image at alpha=1, sites drawn
%       once with a fixed seed and shared by all models (s4/audit_real.py), lung masks from a published network, bone overlap from independent masks.
%   5.3 Why the audit is not circular: the regulariser pastes lesions to teach invariance and the audit pastes lesions to measure it. Mitigations:
%       held-out lesion library, different image domain, independent rib masks, and evidence that does not use insertion (H4' in SAP_v2.md).
%   5.4 Variants: exact-truth audit on the Qure-projection test split (secondary), SZCH audit (development evidence only), amplitude sweep.
%   5.5 Statistics (SAP_v2.md H1', H3'): image-cluster bootstrap, slope of A on log2(1+w), Holm over contrasts.
\subsection{Definition}
For a suppressor $f$ that outputs the bone image, an image $x$ and a lesion signal $n\ge0$ in composite-image units, the absorption at amplitude
$\alpha$ is
\begin{equation}
A=\frac{\big\langle f(x+\alpha n)-f(x),\;\alpha n\big\rangle}{\big\langle \alpha n,\;\alpha n\big\rangle},
\label{eq:absorption}
\end{equation}
with both inner products taken over the lesion support. $A=0$ means that the lesion stays in the soft-tissue image and $A=1$ that it was removed with
the bone; $A$ is a regression slope and can fall outside $[0,1]$. A model that does nothing has $A=0$, so $A$ is reported together with a suppression
gain and never alone. Two suppression measures are paired with $A$: the bone-region PSNR on the held-out SZCH pairs and, on VinDr-RibCXR, the
reduction of gradient energy on rib-boundary pixels (the rib benchmark).

\subsection{Sites and lesions}
Lesions come from the CTs of the LIDC-IDRI test split only, so no audited lesion was seen in training. On each real radiograph six sites are drawn at random
positions inside the lung fields at $\alpha=1$; the sites depend only on the image, its lung mask, the lesion library and a fixed seed, and are shared by
every model. Lung masks come from the PSPNet anatomical segmentation model of TorchXRayVision \citep{cohen2022torchxrayvision}, eroded by
49~px at $1024$~px so that a site lies inside the lung field; the model comes from the ChestX-Det project \citep{lian2021chestxdet}, whose paper does not fully
document its training data. The library holds 296 nodules
(4.0 to 37.6~mm, median 7.1~mm); six sites per image give 1,470 sites on VinDr-RibCXR, 1,482 on JSRT and 450 on SZCH. Because the PSPNet masks
differ slightly between GPUs, evaluations run on cloud GPUs read the lung masks exported on the machine of the final table, so every model is
scored on identical sites.
The bone overlap of a site is measured with independent information, never with a model output: the radiologist rib masks of VinDr-RibCXR
\citep{nguyen2021ribcxr} for the confirmatory audit, and the real dual-energy bone signal on SZCH. A site \emph{overlaps bone} when at least 15\,\% of its
support lies on an annotated rib. On JSRT no bone reference exists; the plan named the consensus of the bone maps of ours and of the reference
model, but the audit took the overlap from the reference model's bone map alone (secondary analysis; recorded as a deviation).

\subsection{Why the audit is not circular}
The twin regulariser pastes lesions to teach invariance, and the audit pastes lesions to measure it, so a model could pass the audit by learning the
insertion artefact. Four design choices limit this. The audit lesions come from CTs that were never used for training; the audit images are real
radiographs, whereas the training composites are synthetic; the confirmatory audit uses rib masks that no model has seen; and the claim that the
regulariser helps is also tested with evidence that does not use insertion: contrast of real nodules on JSRT and Node21 and detection recall of real
lesions grouped by bone overlap (Section~\ref{sec:experiments}). These choices limit the circularity but cannot remove it: a model that recognises
pasted signal in general could still score well, which is why the evidence without insertion is reported beside the audit (Section~\ref{sec:limitations}).

\subsection{Variants}
\emph{Exact-truth audit.} On the Qure-projection test split the lesion is added to the soft component, the true bone image is known, and the
bone error inside the lesion support can be measured against it in addition to $A$. \emph{SZCH audit.} It ran on all 75 SZCH images (450 sites)
during development and is reported as development evidence only (Section~\ref{sec:data}); the SZCH contrasts of
Tables~\ref{tab:audit_contrasts} and~\ref{tab:encoders} use its 51 held-out images (306 sites). \emph{Amplitude sweep.} $A$ is recomputed for
$\alpha\in\{0.25,0.5,1,2,4\}$ on the 450 SZCH sites (development data, Table~\ref{tab:app_amplitude}). For $\alpha\le1$ the absorption is flat
(v1 model 0.301--0.306, no-twin model 0.215--0.218, reference model 0.134--0.136), so the leak is linear in the lesion amplitude; at $\alpha=2$ and
4 it falls for every model, where the added signal starts to saturate at the top of the intensity range.

\subsection{Statistics}
Site-level differences between models are paired; confidence intervals are image-cluster bootstrap intervals. The dose-response over the
twin weight is the least-squares slope $\beta_1$ of
\begin{equation}
\bar A_{s,w}=\beta_0+\beta_1\log_2(1+w)+e_{s,w},
\label{eq:dose}
\end{equation}
where $\bar A_{s,w}$ is the absorption at site $s$ averaged over the seeds trained with twin weight $w$; its interval is an image-cluster
bootstrap interval, and H1$'$ passes if the upper bound is below zero. The contrasts of ours against the reference model and against the
no-twin control are adjusted with the Holm procedure (SAP\_v2.md). Bootstrap intervals use \num{10000} resamples for mean absorption and contrasts,
\num{2000} for the overlap interaction and 500 for the dose-response slope, always resampling whole images with replacement.
